\chapter{Практика 1. Знаковое кодирование}

\section{Задание}
Реализовать знаковое кодирование и декодирование текстового сообщения. Алфавит сообщения состоит из 64 символов: латинские буквы (52 символа), цифры (0...9), пробел и точка. Длина сообщения от 30 до 100 символов. Запрещено использовать встроенные функции Matlab/Octave для преобразования формата данных.

\section{Теоретическое описание}
Знаковое кодирование — это процесс преобразования символьного сообщения в двоичную последовательность. Каждому символу ставится в соответствие уникальное кодовое слово фиксированной длины. Для алфавита из 64 символов достаточно 6 бит ($2^6 = 64$), однако в данной работе используется 8-битное кодирование для удобства.

\section{Реализация}
Кодирование выполняется по следующему правилу:
\begin{itemize}
    \item A-Z: коды 0-25
    \item a-z: коды 26-51
    \item 0-9: коды 52-61
    \item пробел: код 62
    \item точка: код 63
\end{itemize}

Декодер выполняет обратное преобразование: группирует биты по 8, преобразует в число и по таблице символов восстанавливает текст.

\begin{lstlisting}[language=Python, caption=Знаковый кодер и декодер]
class SignCoder:
    BITS = 8

    def sign_encoder(text):
        if not text or len(text) < 30 or len(text) > 100:
            return None
        bits = []
        for c in text:
            if 'A' <= c <= 'Z':
                code = ord(c) - ord('A')
            elif 'a' <= c <= 'z':
                code = ord(c) - ord('a') + 26
            elif '0' <= c <= '9':
                code = ord(c) - ord('0') + 52
            elif c == ' ':
                code = 62
            elif c == '.':
                code = 63
            else:
                return None
            bits.append(format(code, '08b'))
        return ''.join(bits)

    def sign_decoder(bits):
        if not bits or len(bits) % 8 != 0:
            return None
        text = []
        chars = 'ABCDEFGHIJKLMNOPQRSTUVWXYZabcdefghijklmnopqrstuvwxyz0123456789 .'
        for i in range(0, len(bits), 8):
            code = int(bits[i:i + 8], 2)
            if code < 64:
                text.append(chars[code])
            else:
                return None
        return ''.join(text)
\end{lstlisting}

\section{Результат}
Исходное сообщение \textit{"Hello World. This is test message and no more..."} длиной 51 символ было преобразовано в 408 бит. При обратном декодировании сообщение восстановлено полностью без ошибок.

\begin{verbatim}
Исходное сообщение: Hello World. This is test message and no more...
Длина сообщения: 51 символов
Символьное кодирование: 408 бит
Полученное сообщение: Hello World. This is test message and no more...
\end{verbatim}

\section{Вывод}
Знаковый кодер и декодер успешно реализованы. Преобразование текста в битовую последовательность и обратно выполняется корректно.